\section{Conclusion}
\label{sec:conclusion}

We have introduced Phase~2 \cvdgeo{}, the geographic extension of the
Centroidal Voronoi Districts algorithm from B.22~\citep{dellaLibera2026cvd}.
Phase~2 replaces BFS hop-count distance with Euclidean distance on
EPSG:5070 Albers-projected tract centroids, uses the true population-weighted
mean coordinate as the centroid update, and terminates on an
$\varepsilon = 1.0$\,m tolerance to prevent discrete snap oscillation.

\paragraph{When to use geographic CVD.}
The choice between Phase~1 and Phase~2 depends on state geometry:

\begin{itemize}
\item \textbf{Use \texttt{--cvd-metric geographic}} for:
  Florida (panhandle + Keys), Alaska (sparse adjacency, vast tracts),
  Washington (Puget Sound coastal geometry), Maine (peninsulas and island tracts).
  These states have large BFS-to-Euclidean divergence where Phase~2 provides
  6\,\% or more PP improvement.

\item \textbf{Phase~1 (\texttt{graph-distance}) is adequate} for:
  North Carolina, Wisconsin, Iowa, Kansas, and most regular-geometry inland
  states where hop-count and Euclidean distance agree.
  The 2\,\% Phase~2 improvement on NC does not justify the fetch step for
  routine runs.
\end{itemize}

\paragraph{Phase~2 is a strict upgrade for target states.}
Phase~2 improves PP, reduces edge cut, converges faster, and is 13--18\%
faster per run than Phase~1 for the same iteration budget.
There is no quality dimension on which Phase~2 is worse than Phase~1;
the only cost is the one-time centroid fetch step.

\paragraph{YAML configuration.}
The recommended YAML for coastal-state runs:
\begin{verbatim}
algorithm:
  structure: centroidal-voronoi
  weights: geographic
  search: convergence
  cvd_iters: 20
  cvd_metric: geographic
  cvd_seed_init: kmeans++
  convergence_threshold: 600
  balance_tolerance: 0.5
workers: 8
\end{verbatim}

\paragraph{Open questions.}
Several extensions remain for future work:

\begin{enumerate}
\item \textbf{Ellipsoidal Albers.}
  The spherical approximation introduces $\le 3$\,m error per tract.
  An ellipsoidal implementation (PROJ library bindings, GRS80) would
  eliminate this error entirely.
  For redistricting purposes, 3\,m is negligible; for other geographic
  applications of the algorithm, the ellipsoidal formula may be warranted.

\item \textbf{Geographic CVD + BisectionEnsemble.}
  Running \be{}$(T=50)$ with \cvdgeo{} bisectors would select the
  lowest-EC geographic Voronoi plan across 50 base seeds.
  Since Phase~2 is faster per run than Phase~1, the ensemble overhead
  is smaller.
  This configuration is deferred to G.15~\citep{dellaLibera2026comparison}.

\item \textbf{Geographic CVD as SA warm start.}
  Phase~2 plans are geometrically compact and have lower edge cuts than
  Phase~1 plans.
  Using \cvdgeo{} as the initialisation for \sa{} bisection
  (instead of \metis{}) may find lower-EC solutions with different
  compactness properties.

\item \textbf{Ensemble interaction.}
  Verify that \texttt{cvd\_geo\_seed} (which depends on both
  \texttt{base\_seed} and \texttt{node\_path}) produces distinct seeds
  across ensemble members at each bisection node.
  This is expected from the SHA-256 construction but should be asserted
  as an L2 test.

\item \textbf{International territories.}
  Guam and other US territories lack well-established Albers projections.
  A fallback to raw Euclidean on WGS84 coordinates (adequate for
  small territories where meridian convergence is minor) is an option.
\end{enumerate}

\paragraph{Summary.}
Phase~2 \cvdgeo{} completes the geographic generalisation of the \cvd{}
algorithm portfolio.
Phase~1 (B.22) provides a robust, data-light baseline for all states.
Phase~2 (this paper) provides a geometrically exact alternative for states
where geographic and graph-topological distance diverge.
Together they give practitioners a choice: fast and adequate everywhere
(\texttt{graph-distance}), or optimal for coastal and irregular states
(\texttt{geographic}).
